\documentclass[11pt]{article}
\usepackage[a4paper,margin=11mm]{geometry}
\usepackage[no-math]{fontspec}
\setmainfont{Latin Modern Roman}
\usepackage{mathtools,amsmath,amssymb,amsfonts,amsthm,enumerate,arydshln,xurl}
\usepackage[unicode,hidelinks]{hyperref}
\usepackage{thmtools,thm-restate}
\setlength\emergencystretch{3em}
\declaretheorem[name=Theorem, parent=section]{theorem}
\declaretheorem[name=Corollary, sibling=theorem]{corollary}
\declaretheorem[name=Proposition, sibling=theorem]{proposition}
\declaretheorem[name=Lemma, sibling=theorem]{lemma}
\declaretheorem[name=Conjecture, sibling=theorem]{conjecture}
\declaretheorem[name=Observation, sibling=theorem]{observation}
\declaretheorem[name=Fact, sibling=theorem]{fact}
\declaretheorem[name=Definition, sibling=theorem, style=definition]{definition}
\declaretheorem[name=Remark, sibling=theorem, style=definition]{remark}
\declaretheorem[name=Problem, sibling=theorem, style=definition]{problem}
\declaretheorem[name=Example, sibling=theorem, style=definition]{example}
\declaretheorem[name=Claim, sibling=theorem, style=remark]{claim}

\DeclareMathOperator{\Diag}{Diag}
\DeclareMathOperator{\supp}{supp}
\DeclareMathOperator{\subrank}{Q}
\DeclareMathOperator{\slicerank}{SR}
\DeclareMathOperator{\bordersubrank}{\underline{Q}}
\DeclareMathOperator{\tensorrank}{R}
\newcommand{\FF}{\mathbb{F}}
\newcommand{\KK}{\mathbb{K}}
\newcommand{\RR}{\mathbb{R}}
\newcommand{\CC}{\mathbb{C}}
\newcommand{\NN}{\mathbb{N}}
\newcommand{\ZZ}{\mathbb{Z}}
\newcommand{\eps}{\varepsilon}
\DeclareMathOperator{\GL}{GL}

\DeclareMathOperator{\Mat}{Mat}

\DeclareMathOperator{\crk}{crk}
\DeclareMathOperator{\ncrk}{ncrk}
\DeclareMathOperator{\linspan}{span}
\DeclareMathOperator{\rank}{rank} \DeclareMathOperator{\col}{colspan} 

%\DeclareMathOperator{\Diag}{Diag}
%

\DeclareMathOperator{\maxrank}{max-rank}
\DeclareMathOperator{\minrank}{min-rank}
\DeclareMathOperator{\minrnk}{min-rank}
\DeclareMathOperator{\mincov}{min-cov}
\DeclareMathOperator{\minsupp}{min-supp}
\DeclareMathOperator{\maxsupp}{max-supp}


\newcommand{\id}{\mathrm{Id}}
\newcommand{\overCl}[1]{{{#1}^{\overline{\FF}}}}
\newcommand{\overK}[1]{{#1^{\mathbb{K}}}}

% \makeatletter
% \newcommand{\longdash}[1][2em]{%
%   \makebox[#1]{$\m@th\smash-\mkern-7mu\cleaders\hbox{$\mkern-2mu\smash-\mkern-2mu$}\hfill\mkern-7mu\smash-$}}
% \makeatother

%
  \newcommand{\beq}{\begin{equation}}
  \newcommand{\eeq}{\end{equation}}
  \newcommand{\beqn}{\begin{equation*}}
  \newcommand{\eeqn}{\end{equation*}}
  \newcommand{\beqr}{\begin{eqnarray}}
  \newcommand{\eeqr}{\end{eqnarray}}
  \newcommand{\beqrn}{\begin{eqnarray*}}
  \newcommand{\eeqrn}{\end{eqnarray*}}

\IfFontExistsTF{NimbusRoman-Regular}{\newfontface\FixedOriginalUprightFace{NimbusRoman-Regular}}{\PackageError{fixed-source-operators}{Required named NimbusRoman-Regular face missing}{No fallback permitted}}
\makeatletter
% Three finite zero-argument glyphs only. No underaccent or wtilde API is implemented.
\edef\FixedOriginalLowerTilde{\noexpand\mathchar\number\numexpr\symlargesymbols*256+101\relax}
\newcommand{\FixedQGlyphStyle}[2]{\begingroup\setbox0=\hbox{$\m@th#1\text{\normalfont\FixedOriginalUprightFace Q}$}\setbox2=\hbox{$\m@th#1\FixedOriginalLowerTilde$}\dimen0=\dp0\advance\dimen0 by\ht2\advance\dimen0 by1.23pt\typeout{FIXEDGLYPH Q STYLE\meaning#1 BASEwidth=\the\wd0 BASEheight=\the\ht0 BASEdepth=\the\dp0 TILDEwidth=\the\wd2 TILDEheight=\the\ht2 TILDEdepth=\the\dp2 LOWER=\the\dimen0}\hbox{\rlap{\hbox to\wd0{\hss\lower\dimen0\copy2\hss}}\box0}\endgroup}
\newcommand{\FixedRGlyphStyle}[2]{\begingroup\setbox0=\hbox{$\m@th#1\text{\normalfont\FixedOriginalUprightFace R}$}\setbox2=\hbox{$\m@th#1\FixedOriginalLowerTilde$}\dimen0=\dp0\advance\dimen0 by\ht2\advance\dimen0 by1.23pt\typeout{FIXEDGLYPH R STYLE\meaning#1 BASEwidth=\the\wd0 BASEheight=\the\ht0 BASEdepth=\the\dp0 TILDEwidth=\the\wd2 TILDEheight=\the\ht2 TILDEdepth=\the\dp2 LOWER=\the\dimen0}\hbox{\rlap{\hbox to\wd0{\hss\lower\dimen0\copy2\hss}}\box0}\endgroup}
\newcommand{\FixedSRGlyphStyle}[2]{\begingroup\setbox0=\hbox{$\m@th#1\text{\normalfont\FixedOriginalUprightFace SR}$}\setbox2=\hbox{$\m@th#1\FixedOriginalLowerTilde$}\dimen0=\dp0\advance\dimen0 by\ht2\advance\dimen0 by1.23pt\typeout{FIXEDGLYPH SR STYLE\meaning#1 BASEwidth=\the\wd0 BASEheight=\the\ht0 BASEdepth=\the\dp0 TILDEwidth=\the\wd2 TILDEheight=\the\ht2 TILDEdepth=\the\dp2 LOWER=\the\dimen0}\hbox{\rlap{\hbox to\wd0{\hss\lower\dimen0\copy2\hss}}\box0}\endgroup}
\DeclareMathOperator{\asympsubrank}{\mathpalette\FixedQGlyphStyle\relax}
\DeclareMathOperator{\asymprank}{\mathpalette\FixedRGlyphStyle\relax}
\DeclareMathOperator{\asympslicerank}{\mathpalette\FixedSRGlyphStyle\relax}
\makeatother
\begin{document}\begin{center}\Large Problem 2473: Asymptotic subrank of narrow tensors\end{center}\noindent\textbf{Authors:} Jop Briët, Matthias Christandl, Itai Leigh, Amir Shpilka and Jeroen Zuiddam\par\noindent\textbf{Work/version:} Discreteness of Asymptotic Tensor Ranks; exact arXiv2306.01718v4 (9 September 2025); native CC BY4.0; matching Discrete Analysis2025:15, 51pp. (10 September 2025), DOI10.19086/da.143834; printedPDFCC BY3.0\par\noindent\textbf{Native source/license:} \url{https://arxiv.org/src/2306.01718v4}, \url{https://creativecommons.org/licenses/by/4.0/}\par\noindent\textbf{Difficulty:} H2: The proof restricts a bounded-dimensional matrix space to simultaneously diagonal submatrices and converts large maximal rank into a uniform minimum-rank bound. It proves supermultiplicativity and controls mixed tensor powers even when some restricted factors vanish. Chernoff counting supplies exponentially many independent high-minimum-rank slices, and an explicit diagonal-unit extraction turns them into maximal asymptotic subrank.\par\medskip\noindent\textbf{Editorial source boundary:} Theorem 3.17, identical to announced 1.10, is the sole selected narrow-maximal outcome. Complete minimum-rank unit extraction, simultaneous diagonalization, support-removal/basis extension, max-to-min-rank, diagonal supermultiplicativity, mixed zero-factor control and Chernoff/tensor-power closing arguments are retained. Shared tensor/slice definitions remain uncounted context. The cube-root, smallest-dimension-two and discreteness consequences receive no extra count. The author separately announces both general and narrow lower bounds as of independent interest. Exact native v4 CC BY 4.0 and published PDF CC BY 3.0 identities remain separate. Original accents package acquisition is held after HTTP402; no package is cloned. Only the exact fixed zero-argument Q/R/SR under-tilde operators have an approved independent typography resolver using the original installed CMEX0x65 glyph and explicitly named installed NimbusRoman-Regular upright bases. Literal author calls/formulas and source helper definitions remain byte-bound; no arbitrary accent API, unknown argument or font fallback is supported. No mathematical proof, hypothesis or parameter is changed or generated. Upstream classes/logos/archive stay private and unexecuted. Full packet qualification remains pending.\par\medskip\hrule\medskip\noindent\textbf{Original human-authored mathematics follows.}\par
% Exact native author source LF lines 885--978
\setcounter{section}{1}\setcounter{subsection}{3}\setcounter{theorem}{14}
\subsection{Tensor preliminaries}\label{subsec:prelim}

We summarize in this section some standard tensor preliminaries and terminology that we will use throughout the paper, following the papers of Strassen \cite{MR882307,strassen1988asymptotic,strassen1991degeneration} and the book of Bürgisser, Clausen and Shokrollahi~\cite{burgisser1997algebraic}.


\subsubsection*{Tensors}
%
%
Let $\FF$ be any field. For any $n \in \NN$ we denote by $e_1, \ldots, e_n$ the standard basis elements of the vector space $\FF^n$. For any $n_1, n_2, n_3 \in \NN$, we denote by $\FF^{n_1} \otimes \FF^{n_2} \otimes \FF^{n_3}$ the vector space of order-three tensors $T = \sum_{i \in [n_1]} \sum_{j \in [n_2]} \sum_{k \in [n_3]} T_{i,j,k}\, e_i \otimes e_j \otimes e_k$ with coefficients $T_{i,j,k} \in \FF$. In most of this paper we will be dealing only with order-three tensors, so instead of order-three tensor we will just say tensor. 

%
\subsubsection*{Restriction and equivalence} 
We will be using a natural preordering on tensors called restriction, which is defined as follows. For any two tensors $T \in \FF^{n_1} \otimes \FF^{n_2} \otimes \FF^{n_3}$ and $S \in \FF^{m_1} \otimes \FF^{m_2} \otimes \FF^{m_3}$, we say $T$ \emph{restricts to} $S$ and write $T \geq S$ if there are linear maps $L_\ell : \FF^{n_\ell} \to \FF^{m_\ell}$ such that $(L_1 \otimes L_2 \otimes L_3)T = S$. We say $S$ and $T$ are \emph{equivalent} if $T \geq S$ and $S \geq T$. Most tensor properties we will study are invariant under equivalence and monotone under restriction (see \autoref{lem:monotone-restrict}). Finally, we say $S, T \in \FF^{n_1} \otimes \FF^{n_2} \otimes \FF^{n_3}$ are \emph{isomorphic} if $(L_1 \otimes L_2 \otimes L_3)T = S$ for invertible linear maps $L_\ell : \FF^{n_\ell} \to \FF^{n_\ell}$.
%

\subsubsection*{Flattenings and flattening ranks}
Let $T \in V_1 \otimes V_2 \otimes V_3$. Grouping together $V_2$ and $V_3$ gives an element  $T_{1} \in V_1 \otimes (V_2 \otimes V_3)$, and we similarly obtain $T_2$ and $T_3$. %(for any choice of distinct ${\ell_1},{\ell_2},{\ell_3} \in [3]$). 
%We may think of $T_{\ell_1}$ as an $n_{\ell_1} \times (n_{\ell_2}\cdot n_{\ell_3})$ matrix. We call the $T_{\ell_1}$ 
We may think of the $T_\ell$ as matrices and we call them 
the \emph{flattenings} of $T$.
Let $\tensorrank_{\ell}(T)$ ve the matrix rank of $T_{\ell}$. We call $\tensorrank_1(T)$, $\tensorrank_2(T)$ and $\tensorrank_3(T)$ 
%\refnote{7}{missing plural " 's"}\ILnote{added it}\Jshort{I don't like the plural s so much after symbols so I solved it differently} 
the \emph{flattening ranks} of $T$.
%Equivalently, we may slice $T$ into matrices $A_i = (T_{i,j,k})_{j,k} \in \FF^{n_2} \otimes \FF^{n_3}$, $B_j = (T_{i,j,k})_{i,k} \in \FF^{n_1} \otimes \FF^{n_3}$ and $C_k = (T_{i,j,k})_{i,j} \in \FF^{n_1} \otimes \FF^{n_2}$
The triple $(\tensorrank_1(T), \tensorrank_2(T), \tensorrank_3(T))$ is often called the \emph{multilinear rank} of $T$. 


\subsubsection*{Conciseness} 
%\refnote{3}{concise actual definition, out of: \phantomsection\label{ref:3-concise-def13}, \ref{ref:3-first-concise-use},\ref{ref:3-first-concise-def},\ref{ref:3-concise-use7},\ref{ref:3-concise-def8},\ref{ref:3-concise-use10},\ref{ref:3-concise-def13}}
Let $T \in \FF^{n_1} \otimes \FF^{n_2} \otimes \FF^{n_3}$. From the definition of flattening rank it follows that $\tensorrank_\ell(T) \leq n_\ell$ for every $\ell \in [3]$. We call $T$ \emph{concise} if $\tensorrank_\ell(T) = n_\ell$  for every~$\ell \in [3]$. Conciseness essentially says that the tensor cannot fit in a smaller tensor space. The property of being concise is very mild in the following sense. We call a tensor $S$ a \emph{subtensor} of $T = \sum_{i,j,k} T_{i,j,k}\, e_i \otimes e_j \otimes e_k$ if $S$ is of the form $S = \sum_{i\in I,j\in J,k\in K} T_{i,j,k}\, e_i \otimes e_j \otimes e_k$ for some subsets $I,J,K$.

\begin{lemma}\label{lem:equiv-concise}
    Every tensor is equivalent to a concise tensor, which is unique up to equivalence. Concretely, every tensor is equivalent to a subtensor that is concise.
\end{lemma}
\begin{proof}
    The first statement is a standard fact \cite[Section 14.3]{burgisser1997algebraic}. For the second statement, let $T = \sum_{i,j,k} T_{i,j,k}\, e_i \otimes e_j \otimes e_k \in \FF^{n_1} \otimes \FF^{n_2} \otimes \FF^{n_3}$. Let $A_i = (T_{i,j,k})_{j,k} \in \FF^{n_2} \otimes \FF^{n_3}$ for $i \in [n_1]$ be the 1-slices of~$T$. Then one verifies that the dimension of the span of the~$A_i$ equals the flattening rank $r_1=\tensorrank_1(T)$. Suppose $A_1, \ldots, A_r$ are linearly independent, and let $S$ be the tensor with these matrices as 1-slices. Then $S$ is a subtensor of $T$ so clearly $T \geq S$. Also, since the matrices $A_{r+1}, \ldots, A_{n_1}$ are in the span of the matrices $A_1, \ldots, A_r$ we have that $S \geq T$. Since $S$ and $T$ are thus equivalent, their flattening ranks are equal. Repeating this construction for the other two directions gives the claimed subtensor.
    %
    %Zij T een n1 x n2 x n3 tensor. Zij A_1, ..., A_{n1} de 1-slices van T, dus T = \sum_{i=1}^{n1} e_i \otimes A_i. De flattening rank R1(T) is de dimensie van de span van de A_i. Stel A_1 is in de span van A_2, ..., A_{n1}. Laat S = \sum_{i=2}^{n1} e_i \otimes A_i. Dit is een subtensor van T. Er geldt T \geq S en S \geq T. De eerste restrictie T \geq S is duidelijk, want S is een subtensor van T. De tweede restrictie is makkelijk in te zien de slices van S zijn A_2, ..., A_{n1} en A_1 is in de span van die matrices. Met andere woorden S en T zijn equivalent, en dus zijn hun flattening ranks hetzelfde (want flattening ranks zijn monotoon onder restrictie). Met deze procedure kunnen we een tensor U vinden die een subtensor is van T, equivalent is aan T, en concise is.
\end{proof}

Not every space $\FF^{n_1} \otimes \FF^{n_2} \otimes \FF^{n_3}$ contains concise tensors, as can be seen from the following implication (which is in fact known to be an equivalence \cite[Theorem~2]{MR2776439}):

\begin{lemma}\label{cl:concise-formats}
For any $n_1,n_2,n_3\in\NN$, if there is a concise tensor in $\FF^{n_1} \otimes \FF^{n_2} \otimes \FF^{n_3}$, then $n_1 \leq n_2\cdot n_3$, $n_2 \leq n_1\cdot n_3$, and $n_3 \leq n_1\cdot n_2$.
\end{lemma}
\begin{proof}
    Suppose $T \in \FF^{n_1} \otimes \FF^{n_2} \otimes \FF^{n_3}$ is concise. Then $n_{\ell_1} = \rank(T_{\ell_1})$ where $T_{\ell_1}$ is the previously defined flattening of $T$. Since $T_{\ell_1}$ is an $n_{\ell_1} \times (n_{\ell_2}\cdot n_{\ell_3})$ matrix, we have $\tensorrank(T_{\ell_1}) \leq n_{\ell_2} \cdot n_{\ell_3}$. Thus $n_{\ell_1} \leq n_{\ell_2} \cdot n_{\ell_3}$ for every choice of distinct ${\ell_1},{\ell_2},{\ell_3} \in [3]$.
\end{proof}
%
%
%
%

%
%
%
%
%
%
%
%

\subsubsection*{Unit tensors, tensor rank, subrank, and slice rank}
For every $r \in \NN$ we define the tensor $\langle r \rangle = \sum_{i=1}^r e_i \otimes e_i \otimes e_i \in \FF^r \otimes \FF^r \otimes \FF^r$. We call~$\langle r \rangle$ the \emph{unit tensor} of rank $r$.
We say a tensor has \emph{tensor rank one} if it is of the form $u \otimes v \otimes w$ for nonzero vectors $u,v,w$.
The tensor rank of a tensor $T$ is the smallest number $r$ such that $T$ can be written as a sum of $r$ tensors with tensor rank one. 
%
We denote tensor rank by $\tensorrank(T)$. Equivalently, the tensor rank $\tensorrank(T)$ is the smallest number $r$ such that $T \leq \langle r\rangle$. The \emph{subrank} of $T$ is the largest number $r$ such that $\langle r\rangle \leq T$. We denote subrank by $\subrank(T)$. 
%\ILnote{maybe we should add: "For $\ell\in[3]$, we denote by $\subrank_\ell(T)$ the maximal rank of any linear combination of the $\ell$-slices of $T$ (see \autoref{subsec:maxrank}". without it -- we haven't defined it yet but it appears in the lemma below.} 
We say a tensor has \emph{slice rank one} if any of its flattening ranks equals~one. The slice rank of a tensor $T$ is the smallest number $r$ such that $T$ can be written as a sum of $r$ tensors with slice rank one.
%
Subrank, slice rank, the flattening ranks and tensor rank satisfy $\subrank(T) \leq \slicerank(T) \leq \tensorrank_\ell(T) \leq \tensorrank(T)$ for every $\ell\in [3]$.

It is also easy to see that these measures are monotone under restrictions:

\begin{lemma}\label{lem:monotone-restrict}
    For any two tensors $T$ and $S$, if $T\geq S$, then we have $\tensorrank(T)\geq \tensorrank(S)$, $\slicerank(T)\geq \slicerank(S)$, $\subrank(T)\geq \subrank(S)$, and for every $\ell\in[3]$, %$\subrank_\ell(T)\geq \subrank_\ell(S)$ and 
    $\tensorrank_\ell(T)\geq \tensorrank_\ell(S)$.
\end{lemma}


\subsubsection*{Kronecker product and asymptotic tensor parameters}
%
Asymptotic tensor parameters are defined by ``regularizing'' a tensor parameter over large powers under the Kronecker product. For any two tensors 
\begin{align*}
T &= \sum_{i,j,k} T_{i,j,k}\, e_i \otimes e_j \otimes e_k \in \FF^{n_1} \otimes \FF^{n_2} \otimes \FF^{n_3}\\
S &= \sum_{a,b,c} S_{a,b,c} \, e_a\otimes e_b \otimes e_c \in \FF^{m_1} \otimes \FF^{m_2} \otimes \FF^{m_3}
\end{align*}
the Kronecker product $T \boxtimes S \in \FF^{n_1 m_1} \otimes \FF^{n_2 m_2} \otimes \FF^{n_3 m_3}$ is defined as
\[
T \boxtimes S = \sum_{i,j,k} \sum_{a,b,c} T_{i,j,k} S_{a,b,c}\, (e_i \otimes e_{a}) \otimes (e_j \otimes e_{b}) \otimes (e_k \otimes e_{c}).
\]
We define $T^{\boxtimes n}$ as the product of $n$ copies of $T$. The asymptotic tensor rank is defined as $\asymprank(T) = \lim_{n \to \infty} \tensorrank(T^{\boxtimes n})^{1/n}$, which, because of sub-multiplicativity of $\tensorrank$, equals the infimum $\inf_{n} \tensorrank(T^{\boxtimes n})^{1/n}$ (by Fekete's lemma). The asymptotic subrank is defined as $\asympsubrank(T) = \lim_{n \to \infty} \subrank(T^{\boxtimes n})^{1/n}$, which, because of super-multiplicativity of the subrank $\subrank$, equals the supremum $\sup_{n} \subrank(T^{\boxtimes n})^{1/n}$. Over the complex numbers, the asymptotic slice rank is defined as $\asympslicerank(T) = \lim_{n\to\infty} \slicerank(T^{\boxtimes n})^{1/n}$. By the characterization in \cite{MR4495838} this limit exists. Over other fields, we generally do not know whether the limit $\lim_{n\to\infty} \slicerank(T^{\boxtimes n})^{1/n}$ exists, except for oblique tensors. When we do not know whether the limit exists we will instead consider the liminf and refer to this as the ``asymptotic slice rank''.\footnote{The results we have over finite fields for asymptotic slice rank hold equally well using liminf or limsup.}

% Exact native author source LF lines 1019--1098
\setcounter{section}{2}\setcounter{subsection}{0}
\subsection{Max-rank of matrix subspaces and slice spans of a tensor}\label{subsec:maxrank}

%
    %
    We denote by $\FF^{n_1} \otimes \FF^{n_2}$ the space of $n_1 \times n_2$ matrices over $\FF$.
    We define the \emph{max-rank} of any matrix subspace $\mathcal{A} \subseteq \FF^{n_1} \otimes \FF^{n_2}$ as
    %
    \[
\maxrank(\mathcal{A}):=\max\{\rank(A):A\in\mathcal{A}\}.
    \]
    %
    %
    %
    %
%
%
For any tensor $T = \sum_{i,j,k} T_{i,j,k}\, e_i \otimes e_j \otimes e_k \in \FF^{n_1} \otimes \FF^{n_2} \otimes \FF^{n_3}$ we define the matrices %
%
%
%
%
%
%
\begin{align*}
T^{(1)}_i &= \sum_{j,k} T_{i,j,k}\, e_j \otimes e_k \quad (i \in [n_1])\\
T^{(2)}_j &= \sum_{i,k} T_{i,j,k}\, e_i \otimes e_k \quad (j \in [n_2])\\
T^{(3)}_k &= \sum_{i,j} T_{i,j,k}\, e_i \otimes e_j \quad (k \in [n_3]).
\end{align*}
For every $\ell \in [3]$, we call $\{T^{(\ell)}_i : i \in [n_\ell] \}$ the \emph{$\ell$-slices} of $T$. %
%
%
%
%
%
%
%
%
%
%
%
%
%
%
%
%
%
%
%
%
%
%
%
%
%
%
%
%
%
%
%
%
%
%
%
%
%
%
%
%
%
%
We define the $\ell$-slice span of $T$ as the matrix subspace
$\mathcal{A}_\ell = \linspan \{T^{(\ell)}_i : i \in [n_\ell] \}$
and denote its max-rank by
\[
\subrank_\ell(T) = \maxrank(\mathcal{A}_\ell).
\]
%
%
We note that the flattening rank $\tensorrank_\ell(T)$ defined earlier, is equal to the dimension of the subspace $\mathcal{A}_\ell$. 
% Exact native author source LF lines 1938--1940
\setcounter{section}{2}\setcounter{equation}{5}
\section{Min-rank of slice spans and asymptotic subrank of narrow tensors} %
%
\label{sec:lower-bound-imbalanced}
% Exact native author source LF lines 1969--2038

To avoid potential confusion we remark that throughout we will call any matrix $A\in \FF^{n_1} \otimes \FF^{n_2}$ (that is not necessarily square) \emph{diagonal} if its support, which is defined as $\supp(A) := \{(i,j) : A_{i,j} \neq 0\}$, is a subset of $\{(i,i): 1\leq i\leq \min\{n_1, n_2\}\}$.
%

%
%


\subsection{Min-rank of matrix subspaces and its properties}\label{subsec:minrank} %
%
%

%
\newcommand{\StrExC}{T_{\textnormal{NA},c,n}}
%
%
%
%
%
%
%

%
%
%

%

%

%

%
The min-rank of a matrix subspace is simply the smallest rank of any non-zero matrix in the space. 


\begin{definition}
    %
    We define the \emph{min-rank} of a matrix subspace $\mathcal{A} \subseteq \FF^{n_1} \otimes \FF^{n_2}$ as
    %
    \[
    \minrnk(\mathcal{A}):=\min\{\rank(A):A\in\mathcal{A}\setminus\{0\}\}.
    \]
    For any collection of matrices $A_1,
    \ldots,A_c \in \FF^{n_1} \otimes \FF^{n_2}$ we define \[
    \minrnk(A_1,\dots,A_c):=\minrnk(\linspan\{A_1,\dots,A_c\}).
    \]
\end{definition}



%
%
%

%
%
%
%

We will use the following straightforward equivalence. % throughout this section.
%Given matrices $A_1, \ldots, A_c \in \FF^{n_1} \otimes \FF^{n_2}$ we call a linear combination $\alpha_1 A_1 + \cdots + \alpha_c A_c$ with coefficients $\alpha_i \in \FF$ \emph{nontrivial} if at least one of the coefficients $\alpha_i$ is nonzero.

\begin{lemma}\label{lem:minrank-nontrivial-lin-comb}
    Let $A_1,\dots,A_c \in \FF^{n_1} \otimes \FF^{n_2}$ be matrices. %
    For every $r\geq1$, the following two statements are equivalent:
    \begin{enumerate}[\upshape(a)]
        \item The matrices $A_1,\dots,A_c$ are linearly independent and $\minrnk(A_1,\dots,A_c)\geq r$.
        \item For every $u \in \FF^c \setminus \{0\}$ we have $\rank(u_1 A_1 + \cdots + u_c A_c) \geq r$.
    \end{enumerate} 
\end{lemma}
% Exact native author source LF lines 2039--3029

\subsubsection{Large min-rank implies large subrank}\label{subsec:minrank-subrank}

In this section we prove that if the slices of a tensor have large min-rank, then the tensor has large subrank, in the following sense:

%\Jnote{Need to look at the following}

\begin{restatable}{lemma}{lemD}
%
\label{cla-high-rank-high-subrank}
Let $T\in\FF^{n_1}\otimes\FF^{n_2}\otimes\FF^{n_3}$ be a tensor and $A_1,\ldots,A_c\in\FF^{n_1}\otimes\FF^{n_2}$ be $c$ 3-slices of it. %
If $A_1,\ldots,A_c$ are linearly independent and  $\minrnk(A_1,\dots,A_c)\geq 2c(c-1)$,
%\ILnote{the 10 should be $2c(c-1)$} 
then $\subrank(T)\geq c$.    
%
\end{restatable}
%\ILnote{Comment 5. Review A missed the difference between the lemma and the remark. I omitted the remark and changed the lemma to include it already. (the remark is a comment currently)}
%
%
%
%

%
%
%

%
%
%
%

% \begin{remark}
%     It is immediate from \autoref{cla-high-rank-high-subrank} that for any tensor $T\in\FF^{n_1}\otimes\FF^{n_2}\otimes\FF^{n_3}$, if any $c$ linearly independent 3-slices $A_1, \ldots, A_c$ of $T$ satisfy  $\minrnk(A_1, \ldots, A_c)\geq 10c^2$, then $\subrank(T)\geq c$.
% \end{remark}

%

%

%

%
%

%
%
%

\begin{proof}
%We can assume $A_1,\dots,A_c$ are the only $3$-slices (i.e. $n_3=c$) by deleting the other $3$-slices by a restriction.
%
%; some of them will be deleting a row or a column and some of them will be a row or a column operation, i.e. adding a scaled row or column to another in all $3$-slices at once (note that a row operation, for example, is actually adding a scaled $1$-slice to another, meaning performing the same row-operation in all the $3$-slices). 
%We will do this in $c$ steps. 
For $i \in [c]$ and $\alpha \in \{0,1\}$, we say that an $m_1\times m_2$ matrix $A$ has property $P(i,\alpha)$ if $A_{i,i} = \alpha$ and $A_{j, i} = 0$ for all $j \in [m_1]\setminus \{i\}$ and $A_{i,j} = 0$ for all $j \in [m_2] \setminus \{i\}$.

We may assume $n_3 = c$.
We will be changing $T$ with a series of restrictions after which the slices $A_1, \ldots, A_c$ will satisfy: for every $j \in [c]$ the matrix $A_j$ has property $P(j, 1)$ and for every $k \in [c] \setminus \{j\}$ the matrix $A_k$ has property $P(j, 0)$. Then $A_1, \ldots, A_c$ restricted to the first $c$ rows and columns form the unit tensor $\langle c\rangle$.

We will carry out the above inductively, obtaining by induction on $0 \leq i \leq c$, that 
\begin{itemize}
\item $A_1, \ldots, A_c$ are linearly independent and $\minrank(A_1, \ldots, A_c) \geq (c-i)\cdot 2(c-1)$
\item for all $1 \leq j \leq i$, 
\begin{itemize}
    \item $A_j$ has property $P(j, 1)$
    \item for all $k \in [c] \setminus \{j\}$, $A_k$ has property $P(j, 0)$.
\end{itemize}  
\end{itemize}
Throughout the induction, the restrictions we use may change the matrices $A_1, \ldots, A_c$, but we will keep denoting them by these names. In particular, the matrices $A_1, \ldots, A_c$ may become smaller. Namely, in every step of the induction we will remove at most $c-1$ rows and at most $c-1$ columns from all matrices (simultaneously). Note that the case $i=0$ is the assumption of the lemma.

Let $1 \leq i\leq c$ and assume the induction claim holds for $i-1$. 
In the following, all operations on rows and columns we apply to all matrices at the same time.
By linear independence, $A_i$ is nonzero. For $j < i$, row $j$ and column $j$ of $A_i$ are zero (by induction hypothesis), so there are $j,k \geq i$ with $(A_i)_{j,k}$ nonzero.
Swap rows $i$ and $j$ and swap columns~$i$ and $k$ and scale so that $(A_i)_{i,i} = 1$. By adding scalar multiples of row $i$ to rows $j > i$ and adding scalar multiples of column $i$ to columns $k > i$ we make row $i$ and column~$i$ in $A_i$ zero. Then $A_i$ has property $P(i,1)$. 

Now we go over $j \in [c] \setminus \{i\}$ one by one. By adding a scalar multiple of $A_i$ to $A_j$ we get $(A_j)_{i,i} = 0$. Suppose the $i$th row of $A_j$ is nonzero. Then there is a $k>i$ such that $(A_j)_{i,k}$ is nonzero (by induction hypothesis). Add scalar multiples of the $k$th column to the other columns $k'>i$ to make $(A_j)_{i,k'}$ zero. Then remove the $k$th column. Then the $i$th row of $A_j$ is zero. Repeat the same procedure for the $i$th column of $A_j$ to make it zero, possibly removing one row. As a result, $A_j$ has property $P(i,0)$.

For every $j \in [c] \setminus \{i\}$, we removed at most one row and one column, so in total we removed at most $2(c-1)$ rows or columns. Removing a row or column decreases matrix rank by at most one, so (\autoref{lem:minrank-nontrivial-lin-comb})
for any $u\in\FF^{c}\setminus\{0\}$ we have $\rank(u_1A_1+\cdots+u_cA_c)\geq
%$\minrank(A_1, \ldots, A_c) \geq 
(c-(i-1))\cdot 2(c-1) - 2(c-1) = (c-i)\cdot 2(c-1)$. 
%\ILnote{Suggested edits from the email: for the previous sentence: "...by at most one, so for any $u\in\FF^{c}\setminus\{0\}$ we have $\rank(u_1A_1+\cdots+u_cA_c)\geq...$" and for the next sentence: "...are linearly independent and their min-rank is lower bounded by it (\autoref{lem:minrank-nontrivial-lin-comb})."}.
If $i < c$, this is nonzero so $A_1, \ldots, A_c$ are linearly independent and $\minrank(A_1,\ldots,A_c) \geq (c-i)\cdot 2(c-1)$ (\autoref{lem:minrank-nontrivial-lin-comb}). (This is clear for $i=c$ since then the matrices form the tensor $\langle c\rangle$.) %(When $i = c$ linear independence follows since the matrices together from $\langle c\rangle$.) 
This proves the induction claim for $i$.
\end{proof}

% The changes will be split into $c$ steps, each step $i$ making sure the $i$'th $3$-slice $A_i$ has $(A_i)_{i,i}=1$, $0$ in the rest of its $i$'th row and column, and that any other $3$-slice has only $0$ in these row and column. During the $i$'th step we will not be ruining the achievements of previous steps as we will be operating with row and column operation and row and column deletions only for rows and columns $j\geq i$. The previous achievements include having all the at-least-$i$ rows and columns $0$ in their coordinates strictly smaller than $i$ in all $3$-slices (the achievements regarding the rows "protect" the beginning of these columns, for example), so these operations change nothing regarding the rows and columns strictly less than $i$.

% Note that row and column operations, performed for all $3$-slices together, do not change the rank of any matrix in the linear span of them, and thus do not change their min-rank. In each step we will delete at most $c-1$ rows and and at most $c-1$ columns, decreasing the min-rank \ILnote{instead: the rank of any non-trivial linear combination} by at most $2(c-1)$ (as the rank of a matrix is the dimension of its row-span and of its column-span). Hence, at the beginning of each step the min-rank of the $3$-slices is at least $10c^2-2(c-1)c>10c^2-2c^2=8c^2>0$ \ILnote{If I haven't made a mistake this time -- it looks like we can change the 10 into a 2, right?}, meaning they are still linearly independent so no $3$-slice is $0$.

% For the $i$'th step, consider the $i$'th $3$-slice $A_i$, which is nonzero by last observation. By the guarantee from the previous steps we know all the rows and columns $j<i$ are zero, so there is an $j,k\geq i$ with $(A_i)_{j,k}\neq0$. swap between rows $i$ and $j$ and between columns $i$ and $k$ so that $(A_i)_{i,i}=1$ \ILnote{scale}.
% Using row and column operations zero-out the rest of its $i$'th row and column (note again that the ones tha will need to be changed are only with index $j>i$). Use this $3$-slice to zero-out the $(i,i)$'th elements in the other $3$-slices: $(A_j)_{i,i}=0$ for any $j\neq i$.

% Next, for every $j\neq i$, we would like to make the $i$'th column of $A_j$ zero. 
% For this we might need to delete a row: we made sure $(A_j)_{i,i}=0$. If the rest of the column is zero we don't do anything. Otherwise, there is some $k>i$ with $(A_j)_{k,i}\neq 0$. We subtract a scaled version of this row from any other row $k'$ that has $(A_j)_{k',i}\neq0$. Note that these can only be $k'>i$ by the guarantee from previous steps. Finally, we delete row $k$ from our tensor. 

% We follow a similar procedure to zero-out the $i$'th row of $A_j$, with possibly deleting a column from our tensor at the end.

% Note that after these we indeed achieve the required guarantee: we haven't ruined the guarantee from previous steps and we zeroed-out the $i$'th row and column from all $3$-slices apart from $(A_i)_{i,i}=1$. For each $3$-slice except $i$ we deleted at most one row and one column to achieve this, so during the step we indeed deleted at most $c-1$ rows and at most $c-1$ columns.

% When finishing the $c$'th step, we have the tensor $\langle c\rangle$ as the corner subtensor.
% \end{proof}

%
%

%
%
%

%
%

%
%

%
%
%

%
%
%

%
%
%


%
%
%

%
%

%


%

%

%

%

%
%
%
%
%

%
%

%
%



%


%
%
%

%
%
%


%
%
%
%

%
%
%
%

%
%
%

%
%
%

%

%

%

%
%

%
%
%

%
%
%
%

%
%
%
%
%

%

%
%
%

%
%
%
%
%

%
%
%
%
%
%
%

%

\subsubsection{From large max-rank to large min-rank and diagonal}\label{subsec:maxrank-minrank}

In this section we discuss how from any collection of matrices that contains \emph{at least one} matrix of large rank (large max-rank) we can produce a collection of slightly smaller matrices such that \emph{all} nonzero matrices in their span have large rank (large min-rank). Moreover, the resulting matrices are all diagonal, which will be crucial in our application.

For any matrix $A\in \FF^{n_1} \otimes \FF^{n_2}$ %
and subsets $I\subseteq [n_1]$, $J\subseteq [n_2]$ we denote by $A_{I\times J}$ the submatrix of $A$ with rows indexed by $I$ and columns indexed by~$J$.  For any subset $X \subseteq \FF^{n_1} \otimes \FF^{n_2}$ we denote by $X_{I \times J}$ the set $\{A_{I\times J} : A \in X\}$.


\begin{restatable}{lemma}{lemBx}
%\begin{lemma}
\label{lem:minrk_diag}
\label{claim:matrices-into-diagonal-with-large-rank}
For every $c\in\NN$, there is an $\eps(c)>0$ such that the following holds.
%Let $n_1\leq n_2$ be positive integers and 
Let $\mathcal A\subseteq \FF^{n_1\times n_2}$ be a matrix subspace of dimension~$c$.
Then, there exist $U\in \GL(\FF^{n_1})$, $V\in \GL(\FF^{n_2})$, a subset $J\subseteq [\min_\ell n_\ell]$, a basis $B_1,\dots,B_c$ for the space $U\mathcal AV$ and $b\in[c]$ such that:
\begin{itemize}
    \item $(B_1)_{J\times J},\dots,(B_b)_{J\times J}$  are diagonal and linearly independent;
    \item $\minrank\big((B_1)_{J\times J}, \ldots, (B_b)_{J\times J}\big) \geq \eps(c) \maxrank(\mathcal{A})$; % n_2$;
    \item $(B_{b+1})_{J\times J} = \cdots = (B_c)_{J\times J} = 0$.
\end{itemize}
%\end{lemma}
\end{restatable}


Before discussing the proof in detail, as a motivation and ingredient for later, we observe using the pigeonhole principle that any concise tensor has at least one slice of large rank. 

\begin{restatable}{lemma}{lemA}
%
\label{cla:high-rk}
Let $T\in\FF^{n_1}\otimes\FF^{n_2}\otimes\FF^{c}$ be a concise tensor. Then $T$ has at least one 3-slice 
of rank at least $\max\{n_1,n_2\}/c$.    
%
\end{restatable}

\begin{proof}
Suppose $n_1\geq n_2$. %
Concatenate the 3-slices $A_1, \ldots, A_c$ of $T$ to an $n_1\times cn_2$ matrix. Conciseness of $T$ implies that the rank of this matrix is $n_1$ so there are $n_1$ linearly independent columns. By the pigeonhole principle, there is an $i \in [c]$ such that $A_i$ contains at least $n_1/c$ linearly independent columns.   
\end{proof}


% %
% %

% \begin{restatable}{lemma}{lemB}
% %
% \label{claim:matrices-into-diagonal-with-large-rank}
%     Let $0<\eps<1$ and $0<\alpha\leq 1$. For every $c\geq 2$ there is an $N_0(\eps,\alpha,c)\in\NN$ such that for any $n_1,n_2>N_{0}(\eps,\alpha,c)$, and any matrices $A_1,\dots,A_c\in \FF^{n_1}\otimes \FF^{n_2}$ that contain a matrix of rank at least $\alpha \min\{n_1,n_2\}$ in their span, there are invertible matrices 
%     \[
%     R\in \GL_{n_1}(\mathbb{F}),\, C\in \GL_{n_2}(\mathbb{F}),\, S\in \GL_{c}(\mathbb{F})
%     \]
%     and subsets $I\subseteq[n_1],J\subseteq[n_2]$ such that  the matrices in
%     \[
%     X = \biggl\{\Bigl(R\Bigl(\sum_{i=1}^cS_{ij}A_i\Bigr)C\Bigr)\Big|_{I\times J}\biggr\}_{j\in [c]}
%     \]
%     are diagonal, not all zero, have min-rank $\minrnk(X) \geq\min\{n_1,n_2\}^\eps$ and the non-zero elements of $X$ are linearly independent.
% %
% \end{restatable}

% \begin{remark}
%     We stress that in \autoref{claim:matrices-into-diagonal-with-large-rank} it is crucial that we are taking $S$ to be an invertible operation, because we want to maintain the property that the matrices $A_1, \ldots, A_c$ are linearly independent. Even though some matrices will become zero when restricted to $I\times J$ we will use the linear independence of those matrices at a later stage.
% \end{remark}

% %
% %
% The proof of \autoref{claim:matrices-into-diagonal-with-large-rank} has two main parts, which we present separately as the following two lemmas, and which imply it immediately.
% \begin{lemma}\label{lem:matrices-into-diagonal-with-one-large-rank-mat}
%     Let $0<\alpha\leq 1$. For every $c\geq 2$, $n_1,n_2\in \NN$ and any matrices $A_1,\dots,A_c\in \FF^{n_1}\otimes \FF^{n_2}$ that contain a matrix of rank at least $\alpha \min\{n_1,n_2\}$ in their span, there are invertible matrices 
%     \[
%     R\in \GL_{n_1}(\mathbb{F}),\, C\in \GL_{n_2}(\mathbb{F}),\, S\in \GL_{c}(\mathbb{F})
%     \]
%     and subsets $I\subseteq[n_1],J\subseteq[n_2]$ such that  the matrices in
%     \[
%     X = \biggl\{\Bigl(R\Bigl(\sum_{i=1}^cS_{ij}A_i\Bigr)C\Bigr)\Big|_{I\times J}\biggr\}_{j\in [c]}
%     \]
%     are diagonal and one of them is of rank at least $\tfrac{1}{3^{c-1}}\alpha\min\{n_1,n_2\}$.
% \end{lemma}

% \JBnote{I wonder if Lemma~\ref{lem:matrices-into-diagonal-with-one-large-rank-mat} can be proved by induction on~$c$ from the following statement, which I believe is what underlies Algorithm~1: For any $A\in \FF^{n\times n}$ there are invertible matrices $U,V\in \FF^{n\times n}$ and a set $I\subseteq [n]$ of size at least $n/3$ such that the following holds. The principal submatrix $(UAV)_{I\times I}$ is diagonal and for every $x\in \FF^n$, we have that $(U\Diag(x)V)_{I\times I} = \Diag(x)_{I\times I}$.}

% \ILnote{Yes, that's exactly what's going on there. It's a different language for the same thing -- induction instead of iterations (with iteration invariant). If it's more accessible as an induction I have no objection for changing it.}

% %
% %
% %
% \begin{lemma}\label{lem:diag-matrices-with-one-large-rank-mat-into-large-minrank}
%     Let $0<\eps<1$ and $0<\alpha\leq 1$. For any $c\geq2$ there is an $N_0(\eps,\alpha,c)\in\NN$ such that for any $n_1,n_2>N_0(\eps,\alpha,c)$, and any matrices $A_1,\dots,A_c\in\FF^{n_1}\otimes\FF^{n_2}$ and sets of indices $I\subseteq[n_1],J\subseteq[n_2]$ such that $A_1|_{I\times J},\dots,A_c|_{I\times J}$ are diagonal and one of them is of rank at least $\tfrac{1}{3^{c-1}}\alpha\min\{n_1,n_2\}$, there is an invertible matrix $S\in \GL_c(\FF)$ and index sets $\tilde I\subseteq[n_1],\tilde J\subseteq[n_2]$ such that 
%     \[X=\left\{\left(\sum_{i=1}^cS_{ij}A_i\right)\Big|_{\tilde I\times \tilde J}\right\}_{j\in[c]}\]
%     are diagonal, not all zero, have min-rank $\minrnk(X)\geq \min\{n_1,n_2\}^\eps$ and the non-zero elements of $X$ are linearly independent.
% \end{lemma}
% %
% %
% %
% \begin{proof}[Proof of \autoref{lem:matrices-into-diagonal-with-one-large-rank-mat}]
%     Assume $n_1\leq n_2$ (otherwise swap their names).
    
%     Note that the invertible matrix $S$ in the statement of the lemma
%     is just replacing $A_i$'s with linear combinations of them, in a way that keeps their span unchanged, and this is how we will treat it. Similarly, $R$ is invertible row operations done for all the matrices together and $C$ is invertible column operations done for all the matrices together. 
%  %

%  %
%  We start by describing the way $R,C,S$ and $I,J$ are constructed and showing the diagonal property this construction guarantees, together with having a high-rank matrix. 
 
%  %
%  %

%  %
 
%  Without loss of generality (using the $S$'s) we may assume that $\rank(A_1) \geq \alpha n_1$ and that $(A_1)|_{[\alpha n_1]\times [\alpha n_1]} = \id_{\alpha n_1}$ (assume $\alpha n_1\in\NN$. We avoid writing $\lceil\alpha n_1\rceil$ to simplify notation).

% %
 
%  At the beginning, set $I=[\alpha n_1]$. We will iteratively remove indices from it as we go along. $J$ will just be the same set $J=I$. We work in two nested loops: the outer loop (the for-loop in \autoref{alg:diagonalise-large-minrank} below)  goes through the matrices $A_k$ for $k>1$ and makes each diagonal in $I\times I$. It does so by the inner loop (the While loop in \autoref{alg:diagonalise-large-minrank})  which goes over $I$ and changes the matrices by row and column operations and removes some of the indices from $I$. For example, to zero-out a row (except its diagonal element) we use a column which is non-zero in the desired row and subtract scaled versions of that column from all other columns. After that we take the subtracted column's index out of $I$. This process is detailed in the form of an algorithm in \autoref{alg:diagonalise-large-minrank} below.

% \begin{algorithm}[H]
% \caption{Diagonalise with a high-rank element: Construction of $R,C$ and the beginning construction of $I,J$}\label{alg:diagonalise-large-minrank}
%     \begin{algorithmic}
%     \Require $A_1|_{[\alpha n_1]\times[\alpha n_1]}=\id_{\alpha n_1}$
%     \State $R,C\gets\emptyset$ \Comment{$R,C$ will be sets of row and column operations, resp.}
%     \State $I_1\gets[\alpha n_1]$ 
%     \For{$k=2,\dots,c$}
%         \State $I_k\gets I_{k-1}$ \Comment{$I_{k}$ is the current set of indices, some of which we will remove}
%         \State $X_k\gets \emptyset$ \Comment{$X_k$ is the set of indices that we have dealt with}
%         \While{$I_k\setminus X_k\neq\emptyset$}
%             \State $i_k\gets \min I_k\setminus X_k$
%             \If{$\exists j_1\in I_k\setminus\{i_k\}:(A_k)_{i_k,j_1}\neq0$}
%                 \State Subtract a multiple of column $j_1$ from every column $j\in I_k\setminus\{i_k,j_1\}$,
%                 \State making $(A_k)_{i_k,j}$ zero
%                 \State Add these column operations to $C$
%                 \State Remove $j_1$ from $I_k$
%             \EndIf \Comment{Now row $i_k$ in ${A_k}|_{I_k\times I_k}$ is zero outside the diagonal}
%             \If{$\exists j_2\in I_k\setminus\{i_k\}:(A_k)_{j_2,i_k}\neq0$}
%                 \State Subtract a multiple of row $j_2$ from every row $j\in I_k\setminus\{i_k,j_2\}$,
%                 \State making $(A_k)_{j,i_k}$ zero
%                 \State Add these row operations to $R$
%                 \State Remove $j_2$ from $I_k$
%             \EndIf \Comment{Now column $i_k$ in ${A_k}|_{I_k\times I_k}$ is zero outside the diagonal}
%             \State Add $i_k$ to $X_k$
%         \EndWhile
%     \EndFor
%     \State $I\gets I_c$
%     \State $J\gets I$
%     \end{algorithmic}
% \end{algorithm}

% %
 
%  %
% We claim the following iteration invariant: At the end of each iteration of the inner loop (the While loop)
% %
% \begin{enumerate}[(1)]
%     \item  $A_k|_{I_k\times I_k}$'s rows and columns with indices at most $i_k$ are all zero except in their diagonal elements,
%     \item $A_{k'}|_{I_k\times I_k}$ for $k'<k$ are diagonal, 
%     \item and $A_1|_{I_k\times I_k}=\id_{|I_k|}$.
% \end{enumerate}
% %
% %
% We now go over the items of the invariant carefully to make sure they hold.
% The first item ($1$) holds by the conditionals and the operations inside of them -- which explicitly guarantee it, first for the row and then for the column. If we go into the first conditional, say, and subtract a scaled version of column $j_1$ from the other columns, then $(A_k)_{i_k,j_1}\neq0$ but by ($1$) for the previous iteration we know that $j_1>i_k$ (and similarly $j_2>i_k$ in the second conditional), and after this operation the $i_k$'th row of $A_k|_{I_k\times I_k}$ is zero except in its $i_k,j_1$ coordinates. Additionally, by item ($2$) of the invariant in $A_{k'}$ for $k'<k$ the $j_1$ column is nonzero only on the diagonal, so this operation can blemish only the $j_1$'th row of $A_{k'}$ and as we remove $j_1$ from $I_k$ we keep this item as well. Lastly, note that every row or column that may change in $A_1$ have their indices removed from $I_k$ when that happens so indeed we keep $A_1|_{I_k\times I_k}=\id_{|I_k|}$.

% By looking at the end of the final iteration, this invariant shows that when the algorithm terminates we have $A_1|_{I\times J}=\id_{|I|}$ and all of $A_k|_{I\times J}$ are diagonal. 
% We showed above that whenever we go into one of the conditionals the picked $j_1$ or $j_2$ must be $j_1,j_2>i_k$. Thus, every index put into $X_k$ cannot be removed from $I_k$ in a later iteration as a $j_1$ or $j_2$ (as long as we're in the same iteration $k$ of the outer loop). At the end of the iteration for $k$ of the outer loop we have $X_k=I_k$ and each index in that set was responsible to at most $2$ indices removed from $I_k$, implying $|I_k|\geq \frac{1}{3}|I_{k-1}|$. So $|I|=|I_c|\geq\frac{1}{3^{c-1}}|I_1|=\frac{1}{3^{c-1}}\alpha n_1$. Specifically, $\rank(A_1|_{I\times J})\geq\frac{1}{3^{c-1}}\alpha n_1$ by item ($3$) of the invariant.
% \end{proof}
 
%  %
%  %
 
%  %

%  %
%  %
% %

%  \begin{proof}[Proof of \autoref{lem:diag-matrices-with-one-large-rank-mat-into-large-minrank}]
%  We need to guarantee that every linear combination of the nonzero diagonal matrices will be of large rank (which is equivalent to being linearly independent with large min-rank by \autoref{lem:minrank-nontrivial-lin-comb}). 
%  %
%  To achieve this we will replace matrices that appear in a low rank combination with that low rank matrix. We shall then restrict our attention to the complement of the nonzero entries of each such matrix. 

%  Assume w.l.o.g. $n_1\leq n_2$.
%  We know one of the diagonal matrices $A_i$ is of rank at least $\tfrac{1}{3^{c-1}}\alpha n_1$, so we may remove from $I$ and $J$ any index $j$ for which $(A_i)_{j,j}=0$ and we still have all the hypothesis of the lemma met. We may also change their order so that $A_i$ is the first of them (by recording the change of order into the matrix $S$, in which case at the beginning of the algorithm below $S$ is not empty).

%  The process described above is detailed precisely in \autoref{alg:zero-or-highrank} below.

% %
% %

% %
% %
 
% \begin{algorithm}[H]
%     \caption{Make the elements zero or high rank: Construction of $S$ and finishing the construction of $I,J$}\label{alg:zero-or-highrank}
%     \begin{algorithmic}
%         \Require $(A_1)|_{I\times J}=\id$
%         \State $S\gets\emptyset$ \Comment{$S$ will hold instructions to replace a matrix by a} 
%         \State {\hspace{21em}linear combination that includes it}
%         \State $Z\gets\emptyset$
%         \For{$k=1,\dots,c$}
%             \If{$(A_k)|_{I\times J}\neq0$}
%                 \State Add $k$ to $Z$
%             \EndIf
%         \EndFor
%         \While{$\exists A\in\linspan\{A_k\}_{k\in Z}:\rank(A|_{I\times J})<n_1^\eps$}
%             \State Let $A=\sum_{k\in Z}a_k A_k$
%             \State $k_0\gets \max\{k\in Z:a_k\neq0\}$
%             \State $A_{k_0}\gets A$
%             \State Add this replacement to $S$
%             \State Remove the coordinates $i\in I$ for which $A_{i,i}\neq0$ from $I$ and $J$
%             \State Remove $k_0$ from $Z$
%         \EndWhile
%     \end{algorithmic}
% \end{algorithm}

% Note that $Z$ becomes smaller by one in each iteration of the While loop, so the loop has to end in at most $c$ iterations. Additionally, we will show that under some choice for $N_0(\eps,\alpha,c)$, $k_0>1$ in all the iterations, meaning $1\in Z$ at the end of the While loop and thus $Z\neq\emptyset$ at the end of the algorithm, and all matrices in the span of $\{A_k\}_{k\in Z}$ have rank at least $n_1^\eps$ when restricted to rows in $I$ and columns in $J$. The $A_k$'s are still all diagonal in $I\times J$ as taking linear combinations of diagonal matrices keep them diagonal, so all we have left is to explain the choice of $N_0(\eps,\alpha,c)$.

% To guarantee the index $k_0$ in the While loop is always $k_0>1$, we would like to have $|I|\geq n_1^\eps$ at all iterations. That way, $\rank((A_1)|_{I\times J})=\rank(\id_{|I|})\geq n_1^\eps$ and so $A$ with a smaller rank must have at least one matrix $A_k$ with $k>1$ and $a_k\neq0$.

% Recall $|I|\geq\frac{1}{3^{c-1}}\alpha n_1$ at the beginning of this algorithm, and note that in each iteration we remove at most $n_1^\eps$ coordinates from it (the rank of $A$). There are at most $c-1$ iterations before considering $A_1$, so $I$ is of size $|I|\geq\alpha n_1/3^{c-1}-(c-1)n_1^\eps$. As $\alpha n_1/3^{c-1}=\omega(cn_1^{\eps})$ (by the assumptions $\eps<1$ and $\alpha>0$), there is $N_0(\eps,\alpha,c)\in\NN$ such that for any $n_1\geq N_0(\eps,\alpha,c)$ we have $\frac{\alpha}{3^{c-1}}>c n_1^\eps$ and thus $|I|>n_1^\eps$ as required.
%  %
%  %
%  %
% \end{proof}

% %
% %

% \subsubsection{Induction proofs of the above}

\paragraph{Diagonalization.}
The proof of \autoref{lem:minrk_diag} has three parts.
In this first part, we show in \autoref{lem:diagonalize} that any collection of matrices can be made simultaneously diagonal on a large principal submatrix in such a way that the property of having large max-rank is preserved. 
%\Jnote{Added:}
For $a,b \in [n]$ let $E_{a,b} \in \FF^n \otimes \FF^n$ be the matrix with coefficient one at index $(a,b)$ and coefficient zero elsewhere.

\begin{lemma}\label{lem:diagonalize-ind2}
    Let $A \in \FF^{n \times n}$. Then there are $U,V \in \GL(\FF^n)$ and $i,j \in \{2,\ldots, n\}$ such that for $I = [n] \setminus\{i,j\}$,
    %\refnote{16}{the notation for $E_{a,b}$ is only defined much later, on page 39}\Jshort{Fixed above}
    \begin{enumerate}[\upshape (i)]
    \item $(UAV)_{I\times I}  = \alpha_{1,1} E_{1,1} + \sum_{a,b \geq 2} \alpha_{a,b} E_{a,b}$ for some $\alpha_{a,b} \in \FF$;
    \item for every diagonal matrix $B\in \FF^{n\times n}$, we have that $(U B V)_{I\times I} = B_{I\times I}$.
    \end{enumerate}
\end{lemma}
\begin{proof}
    Let $X = \{\ell \in \{2, \ldots, n\} : A_{1, \ell} \neq 0\}$. If $X = \emptyset$, then we let $V = \id$ and let~$j$ be any element of $[n]$. Otherwise, we let $j$ be any element of $X$, and we let $V$ be the matrix that, when multiplied with $A$ from the right, adds multiples of the $j$th column of $A$ to the $\ell$th column of $A$ such that the $(1, \ell)$ entry becomes zero, for every $\ell \in X \setminus \{j\}$%
    %\refnote{17}{$i$ should be $j$}\ILnote{fixed it}
    . In other words, we let
    \[
    V = \id - \sum_{\mathclap{\ell \in X \setminus \{j\}}} A_{1,j}^{-1} A_{1,\ell} E_{j, \ell}.
    \]
    Let $Y = \{\ell \in \{2, \ldots, n\} : A_{\ell, 1} \neq 0\}$. If $Y= \emptyset$ or $Y = \{j\}$, then we let $U = \id$ and we let $i$ be any element of $[n] \setminus \{j\}$. Otherwise, let $i$ be any element of $Y \setminus \{j\}$, and we let $U$ be the matrix that, when multiplied with $A$ from the left, add multiples of the $i$th row of $A$ to the $\ell$th row of $A$ such that the $(\ell, 1)$ entry becomes zero, for every $\ell \in Y \setminus \{i,j\}$.
    In other words,
    \[
    U = \id - \sum_{\mathclap{\ell \in Y \setminus \{i,j\}}} A_{i,1}^{-1} A_{\ell,1} E_{\ell, i}.
    \]
    Then claim (i) is satisfied. 
    To see that (ii) is satisfied we note that $U = \id + U_0$ and $V = \id + V_0$, where~$U_0$ supported by the $i$th column and of~$V_0$ is supported by the $j$th row.
    %Suppose~$B\in \FF^{n\times n}$ is a diagonal matrix.
    Then, for any diagonal matrix $B \in \FF^{n\times n}$,
    \begin{align*}
        UBV &= B + U_0B + BV_0 + U_0BV_0.
    \end{align*}
    Since $i\ne j$, it follows that $U_0BV_0 = 0$.
    Moreover, $U_0B$ is supported by the $i$th column and $BV_0$ is supported by the $j$th row, so $(U_0B)_{I \times I}=0$ and $(BV_0)_{I \times I}=0$, and so $(UBV)_{I \times I} = B_{I\times I}$.
    %the support of $U \Diag(x) V$ is contained in the support of $UV$, which is contained in $\{(i,i) : i \in [n]\} \cup \{(i,\ell): \ell \in [n]\} \cup \{(\ell,j): \ell \in [n]\}$. \Jshort{This is not enough to get (ii) because we wanted to really keep $\Diag(x)$.}
\end{proof}

\begin{lemma}\label{lem:diagonalize-ind}
    Let $A \in \FF^{n \times n}$. Then there are $U,V \in \GL(\FF^n)$ and $I \subseteq [n]$ such that $|I| \geq n/3$ and 
    \begin{enumerate}[\upshape (i)]
    \item $(UA V)_{I \times I}$ is diagonal;
    \item for every diagonal matrix $B \in \FF^{n \times n}$,  we have that $(U B V)_{I \times I} = B_{I \times I}$.
\end{enumerate}
\end{lemma}
\begin{proof}
    The proof is by repeatedly applying \autoref{lem:diagonalize-ind2}.
    Applying \autoref{lem:diagonalize-ind2} to $A$ gives $U,V,I$ such that $(UAV)_{I \times I}$ is of the form 
    \[
    \alpha_{1,1} E_{1,1} + \sum_{a,b\geq 2} \alpha_{a,b} E_{a,b}
    \]
    and such that for any diagonal $B$ we have $(UBV)_{I \times I} = B_{I \times I}$.
    Suppose we are given $A' \in \FF^{m \times m}$ of the form 
    \[
    \sum_{i=1}^k \alpha_{i,i} E_{i,i} + \sum_{a,b\geq k+1} \alpha_{a,b} E_{a,b}.
    \]
    Applying \autoref{lem:diagonalize-ind2} to $A'' = (A')_{\{k+1, \ldots m\} \times \{k+1, \ldots m\}}$ gives $U', V', I'$ such that $(U' A'' V')_{I' \times I'}$ is of the form $\alpha'_{1,1} E_{1,1} + \sum_{a,b\geq 2} \alpha'_{a,b} E_{a,b}$. Then 
    %\refnote{18}{Expecting $\id\oplus U'$ rather than $U'\oplus\id$, based on $A'$'s structure} 
    %\Jshort{I think he is right.}
    \[
    ((\id_k\oplus U') A' (\id_k\oplus V'))_{I' \times I'}
    \]
    is of the form 
    \[
    \sum_{i=1}^{k+1} \alpha_{i,i} E_{i,i} + \sum_{a,b\geq k+2} \alpha_{a,b} E_{a,b}
    \]
    and for any diagonal $B$ we have $((\id_k\oplus U')UBV(\id_k\oplus V'))_{I' \times I'} = B_{I' \times I'}$. In every step of the process $k$ goes up by one, and $|I|$ goes down by two, so that in the end we obtain the claim with $I \geq n/3$.
\end{proof}

\begin{lemma}\label{lem:diagonalize}
Let $A_1 = \id_n, A_2, \ldots, A_c \in \FF^{n \times n}$. Then there are $U,V \in \GL(\FF^n)$ and $I \subseteq [n]$ such that $|I| \geq 3^{-(c-1)} n$ and
\begin{enumerate}[\upshape (i)]
    \item $(U \id_n V)_{I\times I} = (\id_n)_{I\times I}$
    \item $(UA_i V)_{I \times I}$ is diagonal for all $i \in [c]$.
\end{enumerate}
\end{lemma}
\begin{proof}
    The proof is by induction over $c$ using \autoref{lem:diagonalize-ind}. 
    The base case $c = 1$ is clear.
    For the induction step, we are given $A_1, \ldots, A_m$, $U,V$ and $I \subseteq[n]$ with $|I| \geq 3^{-(m-1)}n$ such that for $B_i = (UA_i V)_{I \times I}$ we have $B_1 = \id_I$ and every $B_i$ is diagonal.
    Let $B_{m+1} = (U A_{m+1} V)_{I \times I}$. We apply \autoref{lem:diagonalize-ind} to $A = B_{m+1}$. This gives $U', V' \in \GL(\FF^I)$ and $J \subseteq I$ such that $|J| \geq \tfrac13 |I|$, $(U' B_{m+1} V')_{J \times J}$ is diagonal, and $(U' B_i V')_{J \times J} = (B_i)_{J \times J}$ for every $i \in [m]$. This proves the claim.
\end{proof}

\paragraph{From large max-rank to large min-rank.}
In this second part of the proof of \autoref{lem:minrk_diag}, we show in \autoref{lem:large-minrank} that any low-dimensional matrix subspace~$\mathcal A$ of diagonal matrices with large max-rank contains a principal-matrix subspace with large min-rank.
Moreover, in \autoref{lem:basis-extension-matrices}, we show that any matrix subspace has a basis that, when restricted to a particular principal submatrix, consists of a set of linearly independent matrices and all-zero matrices.

For any vector $w \in \FF^n$ let $\supp(w) = \{i \in [n]: w_i \neq 0\}$ and let $|w| = |\supp(w)|$. For any subspace $W \subseteq \FF^n$ let
\begin{align*}
\maxsupp(W) &= \max \{ |w| : w \in W\},\\
\minsupp(W) &= \min \{ |w| : w \in W \setminus \{0\}\}.
\end{align*}

\begin{lemma}\label{lem:minsupp}
Let $V \subseteq \FF^n$ be a subspace such that $\dim(V) \leq c$.  
%and suppose there is an element in $V$ with support size $s$. 
There is an $I \subseteq [n]$ such that $V_I$ is not zero and
\[
\minsupp(V_I) \geq \tfrac1c \maxsupp(V).
\]
\end{lemma}
\begin{proof}
    We will use the general fact that for arbitrary vectors $w_1, \ldots, w_\ell \in W$, if for every $j\in \{0, \ldots, \ell-1\}$ we have $\cup_{i \leq j} \supp(w_i) \neq \cup_{i \leq j + 1} \supp(w_i)$, then $w_1, \ldots, w_\ell$ are linearly independent.

    Let $k = \maxsupp(V)$ and $d = \dim(V)$.
    Before constructing $I$, we first construct a set $X \subseteq V$ by the following process, starting with $X$ empty.
    If there is a nonzero element $v_1 \in V$ such that $|v_1| < \tfrac1c k$, then add $v_1$ to $X$.
    If $X = \{v_1, \ldots, v_\ell\}$ and there is an element $v_{\ell + 1} \in V$ such that 
    $0 < |\cup_{i \leq \ell+1} \supp(v_i) \setminus \cup_{i \leq \ell} \supp(v_i)| < \tfrac1c k$, then add $v_{\ell+1}$ to $X$.
    By the aforementioned general fact, $X$ is linearly independent and thus $X$ is finite, namely $|X| \leq d$.

    % We claim that in fact $|X| \leq d-1$. Suppose $|X| = d$, then $V$ has a basis $v_1,\ldots, v_d$ satisfying
    % $|\cup_{i \leq \ell+1} \supp(v_i) \setminus \cup_{i \leq \ell} \supp(v_i)| < \tfrac1c k$ for every $\ell$
    % and so we have the inequality $| \cup_{i \in [d]} \supp(v_i)| < d \tfrac1c k \leq k$. However, there is an element $v \in V$ with $|v|  = k$ which cannot be in the span of these $v_1, \ldots, v_d$, which is a contradiction.
    
    We let $J = \cup_i \supp(v_i)$ and $I = [n] \setminus J$.
    
    First we claim that $V_I$ is not zero. We have $|J| < d\tfrac1c k \leq k$.
    Then for any element $v \in V$ with $|v| = k$ we have $|v_I| \geq k - |J| > 1$ and thus $v_I \neq 0$.
    
    Second, we claim $\minsupp(V_I) \geq \tfrac1c k$.  Suppose there is a nonzero element $w \in V_I$ such that $|w| < \tfrac1c k$, then this contradicts the way we constructed $X$.
\end{proof}


\begin{lemma}\label{lem:basis-extension}
    Let $V \subseteq \FF^n$ be a subspace of dimension $c$, and let $I \subseteq [n]$. Then there is a basis $v_1, \ldots, v_c$ for $V$ such that 
    \begin{itemize}
        \item $(v_1)_I, \ldots, (v_b)_I$ are linearly independent
        \item $(v_{b+1})_I = \cdots = (v_c)_I = 0$
    \end{itemize}
    for $b = \dim(V_I)$.
\end{lemma}
\begin{proof}
    Choose any basis $w_1, \ldots, w_c$ for $V$. Let $M$ be the matrix with these basis vectors as columns. The submatrix $M_{I \times[n]}$ has rank $b$ so there is a subset $J \subseteq [n]$ of size $b$ such that $M_{I \times J}$ has rank $b$.
    Let $\{v_1, \ldots, v_b\} = \{w_j : j \in J\}$ and $\{v_{b+1}, \ldots, v_c\} = \{w_j : j \not\in J\}$.
    Then $(v_1)_I, \ldots, (v_b)_I$ form a basis of $V_I$. By subtracting multiples of $v_1, \ldots, v_b$ from the vectors $v_{b+1}, \ldots, v_c$ we can ensure that $(v_{b+1})_I = \cdots = (v_c)_I = 0$.
\end{proof}

\begin{lemma}\label{lem:large-minrank}
    Let $\mathcal{A} \subseteq \FF^{n \times n}$ be a subspace of diagonal matrices such that $\dim(\mathcal{A}) \leq c$. There is an $I \subseteq [n]$ such that 
    \[
    \minrank(\mathcal{A}_{I\times I}) \geq \tfrac1c \maxrank(\mathcal{A}).
    \]
\end{lemma}
\begin{proof}
    This follows from \autoref{lem:minsupp} by identifying every diagonal matrix $A \in \mathcal{A}$ with the vector of its diagonal coefficients $v = (A_{1,1}, \ldots, A_{n,n})$ and $\rank(A) = |\supp(v)|$.
\end{proof}

\begin{lemma}\label{lem:basis-extension-matrices}
    Let $\mathcal{A} \subseteq \FF^{n_1 \times n_2}$ be matrix subspace of dimension $c$, and $J \subseteq [\min_\ell n_\ell]$. Then there is a basis $B_1, \ldots, B_c$ for $\mathcal{A}$ such that
    \begin{itemize}
        \item $(B_1)_{J \times J}, \ldots, (B_b)_{J \times J}$ are linearly independent
        \item $(B_{b+1})_{J \times J}= \cdots = (B_{c})_{J \times J} = 0$
    \end{itemize}
    for $b = \dim(\mathcal{A}_{J\times J})$.
\end{lemma}
\begin{proof}
This follows from applying \autoref{lem:basis-extension} to $V = \mathcal{A}$ and $I = J \times J$.
\end{proof}


\paragraph{Combining the above.} \!We now combine \autoref{lem:diagonalize}, \autoref{lem:large-minrank} and \autoref{lem:basis-extension-matrices} to prove \autoref{lem:minrk_diag}.

\begin{proof}[Proof of \autoref{lem:minrk_diag}]
Let $k = \maxrank(\mathcal A)$.
    By Lemma~\ref{cla:high-rk}, we have that $k \geq n_2/c$.
    After a basis transformation, we may assume that $\mathcal A$ has a basis $A_1 = \id_k,A_2,\dots,A_c$.
    Applying \autoref{lem:diagonalize} to the submatrices $(A_i)_{[k]\times [k]}$, we obtain matrices $U,V\in \GL(\FF^k)$ and a set $I\subseteq[k]$ of size at least $3^{-(c-1)}k$ such that
    \begin{enumerate}[(i)]
        \item $A_1' = \big((U\oplus \id_{n_1 - k})A_1(V\oplus \id_{n_2 - k})\big)_{I\times I} = (\id_k)_{I\times I}$;
    \item $A_i' = \big((U\oplus \id_{n_1 - k})A_i(V\oplus \id_{n_2 - k})\big)_{I\times I}$ is diagonal for all $i\in [c]$.
    \end{enumerate}
    %Let~$V\subseteq \FF^I$ be the vector space spanned by the diagonals of the submatrices~$A_i'$.
    %It follows from~(i) that~$\maxsupp(V) = |I|$.
    %Hence, by Lemma~\ref{lem:minsupp} and the fact that the matrices~$A_i'$ are diagonal that, 
    By \autoref{lem:large-minrank} there is a set~$J\subseteq I$ such that
    \beqn
    \minrnk\big((A_1')_{J\times J},\dots,(A_c')_{J\times J}\big) \geq \tfrac{1}{c} 3^{-(c-1)}k. % \geq \frac{1}{c^2} 3^{-(c-1)} n^2.
    \eeqn
    The claim follows from applying \autoref{lem:basis-extension-matrices} with $B_i = A'_i$.
\end{proof}

\subsubsection{Super-multiplicativity of min-rank}\label{subsec:supermult}

For subspaces of matrices $\mathcal{A}$ and $\mathcal{B}$ their tensor product $\mathcal{A} \otimes \mathcal{B}$ is the subspace of matrices spanned by the products $A \otimes B$ for all $A \in \mathcal{A}$ and $B \in \mathcal{B}$.
In this section we will prove that the min-rank is super-multiplicative under the tensor product, under the assumption that at least one of the matrix subspaces is \emph{diagonal}. (Without this assumption the min-rank can be strictly submultiplicative, see \autoref{ex:minrank-not-supermultiplicative}.)
We then apply this to the min-rank of powers of any diagonal subspace of matrices.
%\ILnote{Comment 12. Changed the parentheses above.}

%
%

\begin{lemma}\label{lem:minrank-diag-supermultiplicative}
    Let $\mathcal{A}$ and $\mathcal{B}$ be subspaces of matrices. Suppose that all matrices in $\mathcal{A}$ are diagonal.
    %
    Then %
$\minrnk(\mathcal{A}\otimes\mathcal{B})\geq\minrnk(\mathcal{A})\cdot\minrnk(\mathcal{B})$.
\end{lemma}
%
\begin{proof}%
    Let $C\in\mathcal{A}\otimes\mathcal{B}$ be a nonzero element. Then $C$ is of the form $C =\sum_{i=1}^\ell A_i\otimes B_i$ for some $A_i \in \mathcal{A}$, $B_i \in \mathcal{B}$ and $\ell \geq 1$. We may assume $\{A_i\}_{i=1}^\ell$ and $\{B_i\}_{i=1}^\ell$ are both linearly independent sets of matrices.  %

    Since every $A_i$ is diagonal, the matrix $C$ is block-diagonal with blocks $D_j$ given by $D_j = \sum_{i=1}^\ell (A_i)_{j,j} B_i$. Thus $\rank(C) = \sum_j \rank(D_j)$. For any $j$ such that $(A_1)_{j,j} \neq 0$, the matrix $D_j$ is a nontrivial linear combination of the $B_i$ and thus $\rank(D_j) \geq \minrnk(\mathcal{B})$. The number of $j$ such that $(A_1)_{j,j} \neq0$ is equal to $\rank(A_1)$ since $A_1$ is diagonal. %
    Thus
    \[
    \rank(C) = \sum_j \rank(D_j) \geq \sum_{\mathclap{j: (A_1)_{j,j} \neq 0}} \rank(D_j) %
    \geq \rank(A_1)\minrnk(\mathcal{B}) \geq \minrnk(\mathcal{A})\minrnk(\mathcal{B}),
    \]
    which proves the claim. %
    %
%
    %
%
    %
\end{proof}
%
%\ILnote{Comment 12. Changed in the example below the "is necessary" to "cannot be removed". We can also phrase it as: "The lemma cannot be made without any condition" if we prefer.}
\begin{example}\label{ex:minrank-not-supermultiplicative}
    The condition in \autoref{lem:minrank-diag-supermultiplicative} that at least one of the matrix subspaces consists of diagonal matrices cannot be removed. To see this we may take the subspace of matrices over the reals
    \[
\mathcal{A} = \Bigl\{\begin{pmatrix}
    a & b\\-b & a
\end{pmatrix} : a,b \in \RR\Bigr\}.
    \]
    Then $\minrnk(\mathcal{A}) = 2$. However, $\mathcal{A}^{\otimes 2}$ contains the element
    \[
    \begin{pmatrix}
        1 & 0\\0 & 1
    \end{pmatrix}^{\otimes 2}
    +
    \begin{pmatrix}
        0 & 1\\-1 & 0
    \end{pmatrix}^{\otimes 2}
     = \begin{pmatrix}
         1 & 0 & 0 & 1\\
         0 & 1 & -1 & 0\\
         0 & -1 & 1 & 0\\
         1 & 0 & 0 & 1
     \end{pmatrix}
    \]
    which has rank $2 < 2^2$, so $\minrnk(\mathcal{A}^{\otimes 2}) < \minrnk(\mathcal{A})^2$.
\end{example}


\begin{lemma}\label{claim:large-rank-lin-comb-tensorprods}
    Let $\mathcal{A}$ %
    be a subspace of matrices. Suppose that all matrices in $\mathcal{A}$ are diagonal. %
    Then for all $k\in\NN$ we have $\minrnk(\mathcal{A}^{\otimes k})\geq \minrnk(\mathcal{A})^k$.
\end{lemma}
\begin{proof}
    Repeatedly apply \autoref{lem:minrank-diag-supermultiplicative}.
\end{proof}
%
%

%

%

%
In the rest of this section we prove a slightly more general version of \autoref{claim:large-rank-lin-comb-tensorprods}.
The reason that we need to do this is that we want to take the output of \autoref{claim:matrices-into-diagonal-with-large-rank} (a collection of linearly independent matrices, such that restricted to specified rows and columns the matrices are diagonal and have large min-rank, or are zero) and use the super-multiplicativity property of \autoref{claim:large-rank-lin-comb-tensorprods} to construct a large collection of matrices (that are tensor products of the original matrices) with large min-rank. However, we cannot do this directly as some of the matrices after restriction are zero. The lemma we prove here deals with that.

%\refnote{19}{The notation $A_{I\times J}$ has been used many times up to this point} 
%\Jshort{Removed it}
%Recall that for any matrix $A \in \FF^{n_1} \otimes \FF^{n_2}$ and subsets $I\subseteq [n_1]$, $J\subseteq [n_2]$ we denote by $A_{I\times J}$ the submatrix of $A$ indexed by rows $I$ and columns $J$. For any subset $X \subseteq \FF^{n_1} \otimes \FF^{n_2}$ we denote by $X_{I \times J}$ the set $\{A_{I\times J} : A \in X\}$.


%
%

%
\begin{restatable}{lemma}{lemC}
%
\label{claim:linear-combinations-of-mixed-tensorprods-high-rank}
    Let $X = \{B_1, \ldots, B_b, C_{b+1} \ldots, C_c\} \subseteq \FF^{n_1} \otimes \FF^{n_2}$ be a set of linearly independent matrices. Let $I\subseteq [n_1]$, $J \subseteq [n_2]$, $|I|=|J|$, such that 
    \begin{itemize}
        \item the matrices $(B_i)_{I\times J}$, for $i = 1, \ldots, b$, are diagonal and linearly independent,
        \item the matrices $(C_i)_{I\times J}$, for $i = b+1, \ldots, c$, are zero.
    \end{itemize}
    For any $m \geq \ell\geq 1$, 
    let $Y$ %
    be the set of all order-$m$ Kronecker products of elements of~$X$ such that at least $\ell$ of the factors are from $B_1, \ldots, B_b$. Then %
    \[
    \minrnk(Y) \geq \minrnk(X_{I\times J})^{\ell}.
    \]
%
\end{restatable}
 
%
%
%
%
%
%

%
\begin{proof}
    %
    %
    Let $A$ be any nontrivial linear combination of elements of $Y$.
    Then
    \[
    A=\sum_{t=1}^r a_t\, D_{i_{t,1}}\otimes\cdots\otimes D_{i_{t,m}}
    \]
    for some nonzero $a_t \in \FF$ and $D_{i_{t, j}} \in X$. Every $D_{i_{t, j}}$ is either in $\{B_1, \ldots, B_b\}$, in which case we will call it a $B$-matrix, or in $\{C_{b+1}, \ldots, C_c\}$, in which case we will call it a $C$-matrix.
    %
    %
    %
    %
    Let $k$ be the maximal number of $B$-matrices that appear as factors in any of the summands. Then $k\geq \ell$. 
    Consider any summand with the maximal number of $B$-matrices in its factors. Suppose it is the first summand. %
    %
    For simplicity of notation, we assume the first $k$ factors of this summand are $B$-matrices and the last $m - k$ factors are $C$-matrices: %
    %
    \[
    D_{i_{1,1}} \otimes \cdots \otimes D_{i_{1,m}}  = B_{i_{1,1}}\otimes\cdots\otimes B_{i_{1,k}}\otimes C_{i_{1,k+1}}\otimes\cdots\otimes C_{i_{1,m}}.
    \]
    We restrict the first $k$ factors of each summand in the definition of $A$ to $I\times J$ to get %\refnote{20}{add "of each summand in the definition"}\ILnote{added it}%
    \[
    A' = \sum_{t=1}^r a_t \, (D_{i_{t, 1}})_{I\times J} \otimes \cdots \otimes (D_{i_{t, k}})_{I\times J} \otimes D_{i_{t, k+1}} \otimes \cdots \otimes D_{i_{t, m}}.
    \]
    If there is a $C$-matrix among $D_{i_{t, 1}}, \ldots, D_{i_{t, k}}$, then $(D_{i_{t, 1}})_{I\times J} \otimes \cdots \otimes (D_{i_{t, k}})_{I\times J}$ is zero. 
    %
    Moreover, nonzero summands cannot have any $B$-matrix among $D_{i_{t, k+1}}, \ldots, D_{i_{t, m}}$, as this would contradict the maximality of $k$. So in any non-zero summand the $D_{i_{t, 1}}, \ldots, D_{i_{t, k}}$ are $B$-matrices and the $D_{i_{t, k+1}}, \ldots, D_{i_{t, m}}$ are $C$-matrices. %So in any non-zero summand has $k$ $B$-matrices and then $m-k$ $C$-matrices.
    %
    %
    %
    We know that the first $k$ factors of the first summand of $A$ are $B$-matrices, so $A'$ is nonzero by the linear-independence of $\{(B_i)_{I\times J}\}^{\otimes k}\otimes\{C_i\}^{\otimes m-k}$.
    %
    %
%
    %
%
    %
%
    Using the notation $B'_i = (B_i)_{I\times J}$, we may rewrite $A'$ as
    %
    % \[
    % A' = \sum_{{j}\in[b]^k}B'_{j_1}\otimes\cdots\otimes B'_{j_k}\otimes\Bigl(\sum_{t=1}^r\delta_{i_{t,1}=j_1}\cdots\delta_{i_{t,k}=j_k}\cdot a_t\, C_{i_{t,k+1}}\otimes\cdots\otimes C_{i_{t,m}}\Bigr).
    % \]
    \[
    A' = \sum_{{j}\in[b]^k}B'_{j_1}\otimes\cdots\otimes B'_{j_k}\otimes M_j,
    \]
    where
    \[
    M_j = \sum_{t=1}^r\delta_{i_{t,1}=j_1}\cdots\delta_{i_{t,k}=j_k}\cdot a_t\, C_{i_{t,k+1}}\otimes\cdots\otimes C_{i_{t,m}}.
    \]
    %We have $a_1\neq0$ so for the multi-index that accompanies it, $(i_{1,1},\dots,i_{1,k})\in[b]^k$, the following is a non-trivial linear combination, meaning a non-zero matrix:
    Note that $S = M_{(i_{1,1},\dots,i_{1,k})}$ is a nonzero matrix, because 
    %Let $\tilde j=(i_{1,1},\dots,i_{1,k})\in[b]^k$ be the first $k$ indices of the first summand in $A$. 
    $a_1\neq0$ and the set $\{C_i\}_{i\in\{b+1,\dots,c\}}^{\otimes(m-k)}$ is linearly independent. %
    % %
    % %
    % \begin{equation}\label{eq:block-lin-comb}
    % M = \sum_{t=1}^r\delta_{i_{t,1}=i_{1,1}}\cdots\delta_{i_{t,k}=i_{1,k}}\cdot a_t\, C_{i_{t,k+1}}\otimes\cdots\otimes C_{i_{t,m}}
    % \end{equation} 
    % is nonzero. 
    %
    %
    %
    %
%
    %
    % We may think of $A'$ as a block matrix composed of $n_1^{m-k}\times n_2^{m-k}$ blocks of size $|I|^{k}\times |J|^{k}$ \refnote{21}{$n_1^{m-k}\times n_2^{m-k}$ block matrix was confusing. Suggested alternative phraseing}\ILnote{changed to the suggestion}. That is, the ``outer'' blocks correspond to linear combinations of products of $C_i$'s, and the ``inner'' blocks correspond to linear combinations of tensor products of $B'_i$'s. 
    %Consider the nonzero linear combination \autoref{eq:block-lin-comb} corresponding to ${\tilde j}$ and denote it by $S$. 
%
    Suppose $S_{x,y}\neq 0$.
    Then
    \[
    \sum_{{j}\in[b]^k}B'_{j_1}\otimes\cdots\otimes B'_{j_k} \cdot (M_j)_{x,y}
    \]
    is a nontrivial linear combination, so it has rank at least 
    \[
    \minrnk(X_{I\times J})^k\geq \minrnk(X_{I\times J})^\ell.
    \]
    This matrix is a submatrix of $A'$, and $A'$ is a submatrix of $A$. This proves the claim.
    %
    % Assume without loss of generality that the $(1,1)$ entry of $S$ is not zero ($S_{11}\neq0$). Block $(1,1)$ in $A'$ is a linear combination of the order-$k$ tensor products of the $B'_i$'s with a nonzero coefficient ($S_{11}$) for $B'_{i_{1,1}}\otimes\cdots\otimes B'_{i_{1,k}}$, meaning this linear combination is nontrivial. \autoref{claim:large-rank-lin-comb-tensorprods} then implies that this block has rank at least $\minrnk(X_{I\times J})^k\geq \minrnk(X_{I\times J})^\ell$ (by $k\geq\ell$) which in turn lower bounds the rank of $A'$ and thus of $A$. %, which proves the claim.
    %\refnote{22}{Page 27, last few lines. These were not very clear to me and I had to read a few times. I hope I now understand, but this step could be expanded slightly.} \Jshort{We rephrased the last part of the proof to make it more concrete.}
\end{proof}

%
%
%
%

%

%
    
%




%
%
    
%
%
    
%
    
%
%
%
%
%
%
%
%
%
%
%
%

%
%
%
%

%
%
%

%
    
%
%
%
%
%
%
%
%
%
%
%
%
%
%

%
%
%
%


\subsection{Narrow tensors have maximal asymptotic subrank}
\label{sec:narrow}\label{subsec:main-unb}



%

%


%
We will now %
%
prove the main theorem of this section, which says that for every constant~$c$, and large enough $n_1, n_2$, all tensors in $\FF^{n_1} \otimes \FF^{n_2} \otimes \FF^{c}$ have maximal asymptotic subrank.

\begin{theorem}%
\label{theorem:concise-cn1n2-assubrank-c}
    For every $c\geq2$ there is $N(c)\in\mathbb{N}$ such that for any $n_1,n_2\in\NN$ such that $\max\{n_1,n_2\}>N(c)$, every concise tensor $T\in\mathbb{F}^{n_1}\otimes\mathbb{F}^{n_2}\otimes\mathbb{F}^c$ satisfies $\asympsubrank(T)=c$.
\end{theorem}

%The proof of \autoref{theorem:concise-cn1n2-assubrank-c} relies on min-rank lemmas that we proved in
%\autoref{subsec:minrank}.
%For convenience of the reader we briefly restate here the ones we need. 
% \refnote{23}{Having the lemmas repeated is indeed convenient -- but it seems like the repetition is excessive}
% \Jnote{I removed them.}

% \lemA*

% \lemBx*

% \lemC*

% \lemD*

%

% \Jnote{Did we fix the Chernoff bound proof below?}
% \ILnote{I didn't try to see how you got to the bound that appears below ($m\geq8c^3$ and $\ell=\tfrac1{2c}m$). We can either do that or revert back to the more loose bound of having $\ell=\tfrac1{100c}m$ and $m>5\ln(4)c$. For the proof of this looser bound see the screenshot from our backup, I upleoded it here as the file "Chernoff\_old.png".}

\begin{proof}%[Proof of \autoref{theorem:concise-cn1n2-assubrank-c}]
Suppose $n_1 \geq n_2$. Suppose $n_1 > N(c)$ for an $N(c)$ to be determined in the proof.
Let $T \in \FF^{n_1} \otimes \FF^{n_2} \otimes \FF^{c}$ be concise.
Let $\mathcal{A} \subseteq \FF^{n_1 \times n_2}$ be the 3-slice span of $T$.
By \autoref{cla:high-rk} we have %there is a slice of $T$ of rank at least $\frac{1}{c} n_1$. In other words, 
$\maxrank(\mathcal{A}) \geq \frac{1}{c} n_1$.
We apply \autoref{lem:minrk_diag} to $\mathcal{A}$, giving matrices $B_1, \ldots, B_b$, $B_{b+1}, \ldots, B_c$ and a subset $J \subseteq [\min_\ell n_\ell]$.
To these, we apply \autoref{claim:linear-combinations-of-mixed-tensorprods-high-rank} for $m \geq \ell \geq 1$ (to be specified later) 
%$m = 8c^3$ and $\ell = \tfrac{1}{2c} m$ 
to obtain a set~$Y$ of slices from~$T^{\boxtimes m}$
%tensor $T'$ with 3-slice span $\mathcal{A}'$ 
satisfying
\beq\label{eq:mrkY}
\minrank(Y) \geq \big(\eps(c)\maxrank(\mathcal{A})\big)^\ell \geq (\eps(c)\tfrac1c n_1)^{\ell}. %\tfrac{1}{100c} m}.
\eeq
The size of~$Y$ equals the number of $m$-tuples in $[c]^m$ containing at least~$\ell$ elements from~$[b]$.
Taking $m \geq 8c$ and $\ell = \tfrac{1}{2c} m$, it then follows from the Chernoff bound (e.g., \cite[Equation~(7)]{MR1045520}) that~$|Y|\geq \tfrac12c^m$.

Since~$T$ is concise, so is~$T^{\boxtimes m}$ and therefore, the set~$Y$ is linearly independent.
Let~$T'$ be the tensor with 3-slices given by the elements in~$Y$.
Then $T^{\boxtimes m} \geq T'$.
Provided
\begin{equation}\label{eq:minrank}
\minrnk(Y) \geq 2(\tfrac12c^{m})^2,
\end{equation}
it follows from \autoref{cla-high-rank-high-subrank} that
\beq\label{eq:QTmbound}
\subrank(T^{\boxtimes m}) \geq \subrank(T') \geq \tfrac12c^m.
\eeq
By our choice of~$\ell$ and~\eqref{eq:mrkY}, we have that~\eqref{eq:minrank} holds if
\beqn
\big(\eps(c)\tfrac1c n_1\big)^{\tfrac{1}{2c}m}
\geq
2(\tfrac12c^{m})^2.
\eeqn
It is clear that there is an~$N(c)\in\NN$ such that this is satisfied for any~$n_1>N(c)$ and every~$m$.
Letting~$m$ go to infinity we then get from~\eqref{eq:QTmbound} that 
\beqn
\asympsubrank(T) = \lim_{m\to \infty}\subrank(T^{\boxtimes m})^{1/m} \geq c.\qedhere
\eeqn
\end{proof}
\begin{thebibliography}{Str88}
\bibitem[BCS97]{burgisser1997algebraic}
Peter B\"{u}rgisser, Michael Clausen, and Mohammad~Amin Shokrollahi.
\newblock {\em Algebraic complexity theory}, volume 315 of {\em Grundlehren der mathematischen Wissenschaften}.
\newblock Springer-Verlag, Berlin, 1997.
\newblock \href {https://doi.org/10.1007/978-3-662-03338-8} {\path{doi:10.1007/978-3-662-03338-8}}.

\bibitem[CK11]{MR2776439}
Enrico Carlini and Johannes Kleppe.
\newblock Ranks derived from multilinear maps.
\newblock {\em J. Pure Appl. Algebra}, 215(8):1999--2004, 2011.
\newblock \href {https://doi.org/10.1016/j.jpaa.2010.11.010} {\path{doi:10.1016/j.jpaa.2010.11.010}}.

\bibitem[CVZ23]{MR4495838}
Matthias Christandl, P\'{e}ter Vrana, and Jeroen Zuiddam.
\newblock Universal points in the asymptotic spectrum of tensors.
\newblock {\em J. Amer. Math. Soc.}, 36(1):31--79, 2023.
\newblock \href {https://doi.org/10.1090/jams/996} {\path{doi:10.1090/jams/996}}.

\bibitem[HR90]{MR1045520}
Torben Hagerup and Christine R\"ub.
\newblock A guided tour of {C}hernoff bounds.
\newblock {\em Inform. Process. Lett.}, 33(6):305--308, 1990.
\newblock \href {https://doi.org/10.1016/0020-0190(90)90214-I} {\path{doi:10.1016/0020-0190(90)90214-I}}.

\bibitem[Str87]{MR882307}
Volker Strassen.
\newblock Relative bilinear complexity and matrix multiplication.
\newblock {\em J. Reine Angew. Math.}, 375/376:406--443, 1987.
\newblock \href {https://doi.org/10.1515/crll.1987.375-376.406} {\path{doi:10.1515/crll.1987.375-376.406}}.

\bibitem[Str88]{strassen1988asymptotic}
Volker Strassen.
\newblock The asymptotic spectrum of tensors.
\newblock {\em J. Reine Angew. Math.}, 384:102--152, 1988.
\newblock \href {https://doi.org/10.1515/crll.1988.384.102} {\path{doi:10.1515/crll.1988.384.102}}.

\bibitem[Str91]{strassen1991degeneration}
Volker Strassen.
\newblock {Degeneration and complexity of bilinear maps: some asymptotic spectra}.
\newblock {\em J. Reine Angew. Math.}, 413:127–180, 1991.
\newblock \href {https://doi.org/10.1515/crll.1991.413.127} {\path{doi:10.1515/crll.1991.413.127}}.

\end{thebibliography}\end{document}
